% !TeX root = ../../../main.tex
\section{Introduction}
\label{sec:intro}

Tissue elasticity has long served as a biomarker for pathology, from manual palpation to modern approaches in ultrasonic elastography~\cite{nightingale_feasibility_2001}. Elastography techniques estimate mechanical properties of tissue by analyzing how tissue moves in response to internal or external forces. One such external force is acoustic radiation force (ARF), where focused ultrasonic pulses are noninvasively generated~\cite{sigrist_Ultrasound_2017,doherty_Acoustic_2013}. ARF-based elastography has become particularly useful for identifying compositional and structural changes in soft organs including the liver~\cite{frulio_Evaluation_2013, barr_Can_2020}, breast~\cite{magalhaes_Diagnostic_2017, krouskop_Elastic_1998,tsai_Sonographic_2013}, thyroid~\cite{zhang_Acoustic_2015}, and prostate~\cite{krouskop_Elastic_1998}.

Two major categories of ARF based methods are acoustic radiation force impulse (ARFI)~\cite{nightingale_feasibility_2001} and shear wave elasticity imaging (SWEI)~\cite{sarvazyan_Shear_1998}. In ARFI, the amplitude of tissue displacement within the ARF excitation region is inversely related to tissue stiffness. However, displacement amplitude is influenced by both tissue stiffness and the ARF amplitude, which is typically unknown, limiting ARFI to qualitative assessments of relative stiffness within a spatial region assumed to experience consistent ARF amplitude. In SWEI, the velocity of ARF induced shear waves is directly related to tissue stiffness. The velocity must be measured over millimeter scale distances, which requires spatial averaging. Additionally, shear wave propagation is affected by wave guidance, out of plane reflections, and geometric dispersion, all of which complicate interpretation in mechanically complex tissues.

To overcome these limitations, we introduce an alternative acoustic radiation force-based method for tissue elasticity quantification called Double Profile Intersection (DoPIo) ultrasound. Instead of relying on shear wave velocity, DoPIo estimates elastic modulus by analyzing the scatterer shearing rate within the region of ARF excitation.
